O HySail é uma proposta de \emph{Proof of Retrievability} orientada a ambientes distribuídos de
armazenamento em nuvem. Sua formulação original busca conciliar três propriedades nor-
malmente difíceis de atender de forma simultânea: disponibilidade do dado, integridade
do conteúdo armazenado e custo reduzido de auditoria [1]. Em vez de exigir a recupe-
ração integral do arquivo a cada verificação, o sistema utiliza meios de recuperação
mais econômicas através de checagens de hash.

Seu funcionamento é estruturado em torno de alguns fundamentos principais:

\begin{enumerate}
\item Upload de partes criptografadas de um arquivo
\item Geração de polinômios no corpo de Galois
\item Elaboração de MACs para cada parte criptografada
\item Checagem da integridade de blocos utilizando os polinômios já apresentados
\item Reconstrução do arquivo original utilizando os blocos necessários e válidos
\item Redistribuição de blocos danificados
\end{enumerate}

Considerando estes pontos, é notável que parte do processo de validação é concentrado
ao lado do usuário, pois, ele é responsável por processos de geração, validação, upload e download.
O servidor apenas armazena as informações, contudo, ainda há uma sobrecarga de responsabilidades
ao lado do usuário, principalmente se existem muitos dados que são distribuidos e validades. Em
vista da possível problemática, o presente trabalho visa propor uma solução que atende os seguintes
critérios:

\begin{enumerate}
\item Diminuir parte das responsabilidade de validação do usuário
\item Garantir que o usuário somente faça o download de arquivos sem que haja muitos
processos intermediários
\item Criar um mecanismo de cobrança que não exija que o usuário faça gastos constantes
com os servidores participantes, e crie um processo de cobrança baseado nas informações
entregues.
\end{enumerate}

Os benefícios existentes em tal aplicação está principalmente baseada em custos de infraestrutura,
semelhante aos utilizados em contextos como AWS em que a cobrança é feita através do uso. Neste
caso, a cobrança será realizada sob os dados armazenados, considerando sua duplicação, e o
processo de pagamento é feito de maneira automática e descentralizada.

\section{Preliminares}

Antes que toda sua arquitetura seja abordada, faz-se necessário explicar alguns dos fundamentos
em torno do tema, de maneira que torna-se simples o desenvolvimento adiante.

\subsection{Ethereum}

O Ethereum é, em geral, descrito como uma plataforma descentralizada de propósito
geral para execução de aplicações distribuídas. Diferentemente de redes voltadas apenas
à transferência de valor, ele também permite registrar estados e regras de negócio por
meio de smart contracts. Como resultado, cria-se um ambiente com maior
transparência, rastreabilidade e menor dependência de intermediários centrais [4, 5].

Em termos funcionais, a rede mantém um registro distribuído no qual contas, saldos e contratos
são atualizados por transações validadas coletivamente. Essas transações são executadas
na Ethereum Virtual Machine (EVM), ambiente que garante o mesmo resultado para a
mesma entrada, independentemente do nó que faz a execução [5]. Com isso, parte da
lógica de confiança sai de serviços centralizados e passa a ser definida em código público
e auditável.

No modelo contemporâneo da rede, a segurança do consenso é sustentada por Proof of
Stake, no qual validadores bloqueiam capital como garantia de comportamento adequado.
Essa mudança reduziu o consumo energético em relação ao modelo anterior e manteve a
propriedade de rastreabilidade pública das operações [6]. Na prática, isso é útil em
cenários que exigem evidências verificáveis, como auditoria, registro de eventos e
liquidação automática de obrigações.

\subsubsection{Smart Contracts}

No contexto do Ethereum, smart contracts são programas imutáveis (salvo padrões de
atualização explicitamente adotados) que ficam na cadeia e executam automaticamente
as condições definidas em código. A ideia de contrato inteligente, em sentido amplo,
aparece em Szabo como um mecanismo computacional para reduzir custos de transação
e ambiguidades na execução de acordos [7]. No Ethereum, esse conceito é implementado
na prática por meio da EVM e de transações que disparam as funções do contrato [5].

Sob a ótica arquitetural desta pesquisa, smart contracts são relevantes por três
razões. Primeiro, permitem explicitar regras de cobrança e distribuição de incentivos
de forma transparente, com menor dependência de reconciliação manual. Segundo,
permitem automatizar pagamentos condicionados a evidências de serviço, o que está
alinhado ao objetivo de vincular remuneração à disponibilidade real dos dados. Terceiro,
favorecem auditoria posterior, pois as decisões de liquidação ficam registradas na rede
[4, 8].

Ainda assim, o uso de smart contracts impõe limitações práticas. Custos de execução,
restrições de desempenho e riscos de implementação (como erros lógicos difíceis de
reverter) exigem cuidado no projeto e na validação do código. Por isso, recomenda-se o
uso de testes, auditorias e mecanismos de contingência. Assim, smart contracts
devem ser entendidos como uma camada complementar de execução verificável, e não
como substituição completa de mecanismos institucionais.

\subsection*{Referências utilizadas}

[4] BUTERIN, V. A Next-Generation Smart Contract and Decentralized Application Platform.
2014. Disponível em: https://ethereum.org/en/whitepaper/.

[5] WOOD, G. Ethereum: A Secure Decentralised Generalised Transaction Ledger (Yellow
Paper). 2014. Disponível em: https://ethereum.github.io/yellowpaper/paper.pdf.

[6] ETHEREUM FOUNDATION. Ethereum Proof-of-Stake (PoS). Disponível em:
https://ethereum.org/en/developers/docs/consensus-mechanisms/pos/. Acesso em: 3 jun. 2026.

[7] SZABO, N. Formalizing and Securing Relationships on Public Networks. First Monday,
v. 2, n. 9, 1997. Disponível em: https://firstmonday.org/ojs/index.php/fm/article/view/548.

[8] ANTONOPOULOS, A. M.; WOOD, G. Mastering Ethereum: Building Smart Contracts and
DApps. Sebastopol: O'Reilly Media, 2018.
